\section{Experimental Setup and Evaluation}
\label{sec:experimental}

\subsection{Experimental Environment}
The ClimateGuardian AI evaluation protocol was executed in a Python 3.12.2 environment, heavily leveraging \texttt{scikit-learn}, \texttt{XGBoost}, \texttt{LightGBM}, and \texttt{PyTorch} for predictive modeling. \texttt{pandas} and \texttt{NumPy} were employed for robust data transformations, while SHAP \cite{Lundberg2017} facilitated interpretability diagnostics. The realtime operational pathway relies on a live interface with the Open-Meteo REST API \cite{OpenMeteo2023}.

\subsection{Model Benchmarking Protocol}
To systematically determine optimal predictive architectures, multiple algorithmic paradigms were evaluated across the three hydro-climatic hazards. The candidate families included Logistic Regression, Random Forest, XGBoost, LightGBM, and Recurrent Neural Networks (LSTM and GRU), along with a Logistic Regression Stacking Ensemble. Within the operational framework, the Random Forest model serves as the primary Flood predictor, while the LightGBM classifier evaluating exactly 41 proxy-remediated features serves as the Heatwave predictor.

For the Compound hazard target, four primary candidates were benchmarked: XGBoost, LightGBM, a Stacking Ensemble, and a temporal GRU. The Stacking Ensemble operates on the meta-probabilities generated by XGBoost and LightGBM bases, mapping them through a Logistic Regression meta-learner. Conversely, the GRU requires a sequential 14-day history formatted as a temporal tensor $(N, 14, 45)$ and functions strictly as an offline temporal benchmark. 

\subsection{Chronological Experimental Protocol and Tuning}
All final model evaluation reflects the chronological framework approximating future operational deployment. Training strictly occurred on the 2000--2012 partition, hyperparameter optimization and generalized calibration operated on the 2013--2015 validation set, and the 2016--2018 testing partition was held out exclusively for final, unseen performance quantification. Where cross-validation was required during tuning, a \texttt{TimeSeriesSplit} was strictly employed to prevent temporal inversion.

Final operational configurations were identified through rigorous hyperparameter tuning:
\begin{itemize}
    \item \textbf{Random Forest:} 100 estimators, maximum depth of 5, \texttt{class\_weight='balanced'}.
    \item \textbf{LightGBM:} learning rate of 0.1, unlimited maximum depth, and \texttt{scale\_pos\_weight} handling class imbalance.
    \item \textbf{GRU:} sequence length of 14, hidden dimension of 32, 0.2 dropout, and Adam optimization (learning rate 0.01). Models were trained utilizing \texttt{BCEWithLogitsLoss} and a \texttt{pos\_weight} tensor for maximum 15 epochs, governed by early stopping patience of 5.
    \item \textbf{Stacking Ensemble:} A Logistic Regression meta-learner ($C=1.0$, balanced class weight) operating on two intermediate input probabilities.
\end{itemize}

\subsection{Evaluation Protocol and Metrics}
Prediction performance is characterized using a diverse portfolio of metrics: Accuracy, Precision, Recall, $F_1$-score, Area Under the Receiver Operating Characteristic Curve (ROC-AUC), Precision-Recall Curve Area (PR-AUC), and the Brier Score. Relying exclusively on Accuracy can be profoundly misleading under severe environmental event imbalance. Therefore, Precision quantifies the operational false-alarm burden, while Recall assesses missed-event sensitivity. The $F_1$-score balances both metrics. ROC-AUC and PR-AUC characterize threshold-independent discrimination, with PR-AUC proving particularly informative under class imbalance. The Brier Score quantifies the absolute calibration of raw predictive probabilities.

\subsection{Comparative and Diagnostic Experiments}
The final predictive performance established across the chronological evaluation is detailed in Table~\ref{tab:model_performance}. To deeply analyze the behavioral mechanisms of the proposed algorithms, the framework implements four specialized experimental pathways:
\begin{itemize}
    \item \textbf{Lead-Time Analysis:} Evaluates temporal prediction degradation by artificially shifting future hazard targets. For the Heatwave hazard (Fig. 6), model performance declines from $F_1 = 0.8500$ at a 1-day forecast horizon to $F_1 = 0.7400$, $0.6500$, and $0.6200$ at 3, 5, and 7-day horizons respectively. 
    \item \textbf{Ablation Study:} To determine optimal feature groupings, the Compound model is iteratively evaluated on independent subspaces (Fig. 7). Configuration $F_1$ scores highlight critical synergies: Climate Only ($0.2400$), Hydro Only ($0.2500$), Climate + Hydro ($0.2286$), Climate + Hydro + Temporal/Lag ($0.4051$), and Full Features ($0.3544$).
    \item \textbf{Calibration:} Investigates whether classification boundary discrimination correlates with reliable probability intervals (Fig. 4). Using the Compound model, uncalibrated Brier Scores ($0.2013$) are contrasted against probabilities corrected by Platt Scaling ($0.0338$) \cite{Platt1999} and Isotonic Regression ($0.0386$) \cite{Zadrozny2002}. 
    \item \textbf{SHAP Interpretability:} Quantifies relative feature attributions independent of linear causality (Fig. 5). The proxy variables strictly excluded from the Heatwave model are concurrently removed from SHAP explanation pipelines to guarantee interpretive purity.
\end{itemize}

\begin{table}[htbp]
\caption{Final Offline Test Performance (2016--2018)}
\begin{center}
\resizebox{\columnwidth}{!}{
\begin{tabular}{l l r r r r r r}
\toprule
\textbf{Hazard} & \textbf{Model} & \textbf{Acc} & \textbf{Prec} & \textbf{Rec} & \textbf{F1} & \textbf{ROC} & \textbf{PR}\\
\midrule
Flood & Random Forest & 0.6916 & 0.3539 & 0.6930 & 0.4686 & 0.7184 & 0.3500 \\
Heatwave & LightGBM & 0.9799 & 0.8788 & 0.8969 & 0.8878 & 0.9969 & 0.9712 \\
Compound & Stacking (Op) & 0.9124 & 0.2821 & 0.7333 & 0.4237 & 0.9434 & 0.3318 \\
Compound & GRU (Offline) & 0.8996 & 0.2876 & 0.9778 & 0.4444 & 0.9629 & 0.3730 \\
\bottomrule
\end{tabular}
}
\label{tab:model_performance}
\end{center}
\end{table}

\subsection{Operational Model Selection}
Table~\ref{tab:op_selection} defines the deployment criteria connecting historical evaluation with live inference constraints. For isolated hazards, the best-performing evaluated model seamlessly aligns with operational deployment. However, for Compound predictions, the Logistic Regression Stacking Ensemble ($F_1 = 0.4237$, ROC-AUC = $0.9434$) provides the strongest offline benchmark in the current evidence, outperforming the GRU ($F_1 = 0.4444$). While the GRU requires a historical sequence input defined by $(N, 14, 45)$, the ClimateGuardian AI operational predictor is deliberately designed around stateless point-in-time feature vectors. Therefore, the Stacking Ensemble is explicitly retained as the operational Compound predictor, representing both absolute offline predictive supremacy and pragmatic deployment-oriented model selection.

\begin{table}[htbp]
\caption{Operational Model Selection Rationale}
\begin{center}
\begin{tabular}{p{1.3cm} p{2.2cm} p{1.6cm} p{2.4cm}}
\toprule
\textbf{Hazard} & \textbf{Offline Benchmark} & \textbf{Operational Model} & \textbf{Operational Reason} \\
\midrule
Flood & Random Forest & Random Forest & Stateless point-in-time evaluation \\
Heatwave & LightGBM & LightGBM & Proxy-leakage remediated \\
Compound & Stacking Ensemble & Stacking Ensemble & Highest performance and stateless point-in-time inference compatibility \\
\bottomrule
\end{tabular}
\label{tab:op_selection}
\end{center}
\end{table}
